\begin{frame}{기본 3종: Sigmoid / Tanh / ReLU}
  \small
  \begin{tabular}{@{}p{1.1cm}p{3.0cm}p{3.4cm}p{4.6cm}@{}}
    \toprule
    함수 & 정의 & 미분 & 특징 \\
    \midrule
    Sigmoid & $\dfrac{1}{1+e^{-z}}$ & 최댓값 0.25 ($z=0$) & 출력 $(0,1)$, 확률
      해석 유용하나 \textbf{소실에 취약} \\
    Tanh & $\dfrac{e^z - e^{-z}}{e^z + e^{-z}}$ & 최댓값 1 & 출력 $(-1,1)$,
      sigmoid보다 낫지만 \textbf{양 끝에서 포화} \\
    ReLU & $\max(0, z)$ & $z>0$: 1, \; $z<0$: 0 & 양수 구간 미분 \textbf{항상
      1} — 층을 거쳐도 작아지지 않음 \\
    \bottomrule
  \end{tabular}
  \vskip0.5em
  \begin{columns}[T]
    \begin{column}{0.62\textwidth}
      \centering
      \includegraphics[width=\linewidth]{ch09_2_activations.png}
    \end{column}
    \begin{column}{0.36\textwidth}
      {\footnotesize 여섯 가지 활성함수(윗줄 sigmoid/tanh/ReLU,
      아랫줄 Leaky ReLU/GELU/Swish)와 각 미분 — sigmoid 미분의 최댓값
      0.25(점선)과 ReLU 미분의 0/1 양극화가 대비.}
    \end{column}
  \end{columns}
\end{frame}
